\hspace{3mm}

\centerline{\bf \huge 10 класс}

\begin{flushright}
        \scriptsize{}
\end{flushright}

\problem{\textbf{1}}
В ряд стоят четыре человека, среди которых есть Аюка (который всегда говорит правду), Басан (который всегда лжёт) и преподаватели Расул Владимирович с Анджуром Аркадьевичем (которые могут сказать как правду, так и солгать).

--- Самый левый сказал: <<Рядом со мной стоит Аюка>>.

--- Второй слева сказал: <<Рядом со мной стоит Басан>>.

--- Третий слева сказал: <<Рядом со мной стоят преподаватели>>.

--- Самый правый сказал: <<Рядом со мной стоит Аюка>>.

Определите, кто и где может стоять.

%https://math.mosolymp.ru/upload/files/2019/khamovniki/9/2019-01-14-olympiad-solutions.pdf

\problem{\textbf{2}}
Даны четыре положительных вещественных числа $a, b, c, d$. Величины $a-c, b-d, c+d$ являются сторонами невырожденного треугольника. Докажите, что $a b \neq 2 c d$.

%Решение. Действуем от противного. Тогда из $a-c>0$ получаем $a>c$, а значит из $a b=2 c d$ получаем, что $b<2 d$. Аналогично $a<2 c$. Складывая, имеем $a+b<2(c+d)$. С другой стороны, из неравенства треугольника $(a-c)+(b-d)>c+d$, откуда $a+b>2(c+d)$, противоречие.

\problem{\textbf{3}}
Лягушонок по имени Баир сидит в точке $(0,0)$ координатной плоскости. Он умеет прыгать на расстояние 5, но только в точки с целыми координатами (то есть на первый прыжок у него есть 12 возможных вариантов). Какое наименьшее количество прыжков ему нужно сделать, чтобы достичь точки $(2026, 2026)$?

\problem{\textbf{4}}
Спортсменка Алина и сонная Яна играют в игру: они по очереди убирают камни из кучи, в которой $n$ камней. Количество камней, убираемых за один ход, должно быть на 1 меньше, чем некоторое простое число (необязательно каждый ход одно и то же). Побеждает игрок, который забирает последний камень. Алина ходит первой. Докажите, что существует бесконечно много значений $n$, для которых у Яны есть выигрышная стратегия (Например, если $n=17$, то у Яны нет выигрышной стратегии).

\problem{\textbf{5}}
На биссектрисе угла $A$ треугольника $ABC$ отмечена точка $D$. Прямые $BD$ и $CD$ пересекают отрезки $AC$ и $AB$ в точках $E$ и $F$ соответственно. Прямая $EF$ пересекает описанную окружность в точках $P$ и $Q$. Докажите, что $DO \perp BC$, где $O$ – центр описанной окружности треугольника $DPQ$.   